% Imported from LaTeX_and_pdfs/DGDT_dump.tex; see provenance.json.
\part{Morse Theory}


\section{Morse Theory introduction from a physicist}

I needed some physical motivation to understand Morse theory, and so I looked at Hori, et. al. \cite{dgdt:Hori2003}.  

cf. pp. 43, Sec. 3.4 Morse Theory, from Ch. 3. Differential and Algebraic Topology of Hori, et. al. \cite{dgdt:Hori2003}.

Consider smooth $f:M \to \mathbb{R}$, with non-degenerate critical points.

If no critical values of $f$ between $a$ and $b$ ($a<b$), then subspace on which $f$ takes values less than $a$ is deformation retract of subspace where $f$ less than $b$, i.e.
\[
\lbrace x \in M | f(x) < b\rbrace \times [0,1] \xrightarrow{ F } \lbrace x \in M | f(x) < b\rbrace
\]
$\forall \, x \in M$ s.t. $f(x) < b$,
\[
\begin{aligned}
  & F(x,0)  = x \\
  & F(x,1) \in \lbrace x \in M | f(x) < a \rbrace 
  \end{aligned} \qquad \qquad \, \text{ and } F(a',1) = a' \qquad \, \forall \, a' \in M \text{ s.t. } f(a') < a
\]

To show this, consider $-\nabla f/|\nabla f|^2$

Morse lemma: $\forall \, $ critical pt. $p$ s.t. $\exists \, $ choice of coordinates s.t.
\begin{equation}
  f  = - (x_1^2 + x_2^2 + \dots + x_{\mu}^2) + x_{\mu + 1}^2 + \dots + x_n^2
\end{equation}
where $f(p)=0$ and $p$ is at origin of these coordinates.

\begin{itemize}
\item difference between
  \[
f^{-1}(\lbrace x \leq -\epsilon \rbrace) , \, f^{-1}(\lbrace x \leq + \epsilon \rbrace)
\]
can be determined by local analysis and only depends on $\mu$, $\mu \equiv $ ``Morse index'' $=$ number of negative eigenvalues of Hessian of $f$ at critical pt.

Answer: \\

\[
f^{-1}(\lbrace x \leq + \epsilon \rbrace) \text{ can be obtained from } f^{-1}(\lbrace x \leq -\epsilon \rbrace) \text{ by ``attaching $\mu$-cell'' along boundary $f^{-1}(0)$ }
\]

\item ``attaching $\mu$-cell to $X$ mean, take
  $\mu$-ball $B_{\mu} = \lbrace |x| \leq 1 \rbrace$ in $\mu$-dim. space, \\
  identity pts. on boundary $S^{\mu-1}$ with pts. in the space $X$, through \\
  cont. $f : S^{\mu-1} \to X$, i.e. take
  \[
X \coprod B_{\mu}
\]
with $x\sim f(x)$ \, $\forall \, x \in \partial B_{\mu} = S^{\mu -1}$.

 
\item find homology of $M$,
  
  $f$ defines chain complex $C_f^*$, $k$th graded piece $C^{\alpha_k}$, $\alpha_k$ is number of critical pts. with index $k$.
  \begin{equation}
\begin{aligned}
  & \partial : C_p^k \to C^{k-1}_p \\ 
  & \partial x_a = \sum_b \Delta_{a,b} x_b 
  \end{aligned}
    \end{equation}
  where $\Delta_{a,b} :=$ signed number of lines of gradient flow from $x_a$ to $x_b$, $b$ labels pts. of index $k-1$.

  Gradient flow line is path $x(t)$ s.t. $\dot{x} = \nabla (f)$, with $\begin{aligned} & \quad  \\
    & x(-\infty) = x_a \\
    & x(+\infty) = x_b \end{aligned}$

\item  To define this number ($\Delta_{a,b} $?), construct moduli space of such lines of flow (???) \\
  by intersecting outward and inward flowing path spaces from each critical point, and then show this moduli space is oriented, 0-dim. manifold (pts. with signs)

\item $\partial^2=0$ proof

  $\partial$, boundary of space of paths connecting critical points, whose index differs by $2 =$ union over compositions of paths between critical pts. whose index differs by $1$.

  $\Longrightarrow$ coefficients of $\partial^2$ are sums of signs of pts. in $0$-dim. space, which is boundary of $1$-dim. space.

  These signs must therefore add to $0$, so $\partial^2=0$.  
  
  \end{itemize}

Hori, et. al. \cite{dgdt:Hori2003} is good for physics, but there isn't much thorough, step-by-step explanations of the math.  I will look at Hirsch (1997) \cite{dgdt:MHirsch1997} and Shastri (2011) \cite{dgdt:AShastri2011} at the same time.  

\subsection{Introduction, definitions of Morse Functions, for Morse Theory}

cf. Ch. 6, Morse Theory of Hirsch (1997) \cite{dgdt:MHirsch1997}, Section 1. Morse Functions, pp. 143-

Recall for $TM$, $T_xM \xrightarrow{\varphi}\mathbb{R}^n$.  \\
Cotangent bundle $T^*M$ defined likewise:
\[
T^*_xM \xrightarrow{ \varphi } \text{ dual vector space } (\mathbb{R}^n)^* = L(\mathbb{R}^n,\mathbb{R})
\]
i.e.
\[
T^*M = \bigcup_{x\in M} (M_x^*) \qquad \qquad \, M_x^* = L(M_x,\mathbb{R})
\]
If chart $(\varphi, U)$ on $M$, natural chart on $T^*M$ is
\[
\begin{aligned}
  & T^*U \to \varphi(U) \times (\mathbb{R}^n)^* \\ 
  & \lambda \in M_x^* \mapsto (\varphi(x), \lambda \varphi_x^{-1} )
  \end{aligned}
\]

Projection map
\[
\begin{aligned}
  & p : T^* \to M \\ 
  & M_x^* \mapsto x
  \end{aligned}
\]
Let $C^{r+1}$ map, $1\leq r \leq \omega$, $f:M \to \mathbb{R}$, $\forall \, x \in M$, linear map $T_x f :M_x \to \mathbb{R}$ belongs to $M_x^*$
\[
T_xf = Df_x \in M_x^*
\]
Then
\[
\begin{aligned}
  & Df:M \to T^*M \\ 
  &  x\mapsto Df_x = Df(x)
  \end{aligned}
\]
is $C^r$ section of $T^*M$.  

\begin{definition}
  \textbf{critical point} $x$ of $f$ is zero of $Df$, i.e. \begin{equation}
Df(x) = 0 
  \end{equation} of vector space $M_x^*$.  
  \end{definition}

Thus, set of critical pts. of $f$ is counter-image of submanifold $Z^* \subset T^*M$ of zeros.  \\
Note $Z^*\approx M$, codim. of $Z^*$ is $n=\text{dim}M$.  

\begin{definition}
  \textbf{Morse function} $f$ if $\forall \, $ critical pts. of $f$ are nondegenerate.
  \end{definition}

Note set of critical pts. closed discrete subset of $M$.  \\
Let open $U \subset \mathbb{R}^n$, let $C^2$ map $g:U\to \mathbb{R}$, \\
critical pt. $p\in U$ nondegenerate iff
\begin{itemize}
\item linear $D(Dg)(p):\mathbb{R}^n \to (\mathbb{R}^n)^*$ bijective
\item identify $L(\mathbb{R}^n, (\mathbb{R}^n)^*)$ with space of bilinear maps $\mathbb{R}^n \times \mathbb{R}^n \to \mathbb{R}$, $\Longrightarrow$ equivalent to condition that symmetric bilinear $D^2g(p) : \mathbb{R}^n \times \mathbb{R}^n \to \mathbb{R}$ non-degenerate 
\item $n\times n$ \emph{Hessian matrix}
  \[
\left[ \frac{ \partial^2 g}{ \partial x^i  \partial x^j }(p) \right]
\]
has rank $n$
\end{itemize}

Hessian of $g$ at critical pt. $p$ is quadratic form $H_pf$ associated to bilinear form $D^2g(p)$
\[
\Longrightarrow H_pf(y) =D^2g(p)(y,y) = \sum_{i,j} \frac{ \partial^2g}{ \partial x^i \partial x^j}(p)y^i y^j
\]

Let open $V \subset \mathbb{R}^n$, suppose $C^2$ diffeomorphism $h: V\to U$.

Let $q=h^{-1}(p)$, so $q$ is critical pt. of $gh:V\to \mathbb{R}$.

\begin{tikzpicture}
  \matrix (m) [matrix of math nodes, row sep=3.8em, column sep=4.8em, minimum width=2.2em]
  {
    \mathbb{R}^n & \mathbb{R} \\ 
     \mathbb{R}^n & \\ 
};
  \path[->]
  (m-1-1) edge node [above] {$H_q(gh) $} (m-1-2)
          edge node [auto] {$Dh(q) $} (m-2-1)
  (m-2-1) edge node [below] {$H_pg $} (m-1-2)
  ;
\end{tikzpicture}  

(quadratic) form $(H_pf)$ invariant under diffeomorphisms.


Let $C^2$ $f:M\to \mathbb{R}$.  \\
$\forall \,$ critical pt. $x$ of $f$, define \\
Hessian quadratic form
\[
\begin{aligned}
  & H_xf : M_x \to \mathbb{R} \\ 
  & H_xf : M_x \xrightarrow{ D\varphi_x } \mathbb{R}^n \xrightarrow{ H_{\varphi(x)}(f\varphi^{-1} ) } \mathbb{R}
\end{aligned}
\]
where $\varphi$ is any chart at $x$.

Thus, critical pt. of a $C^2$ real-valued function nondegenerate iff associated Hessian quadratic form is nondegenerate.  

Let $Q$ nondegenerate quadratic form on vector space $E$.

$Q$ negative definite on subspace $F \subset E$ if $Q(x)<0$ whenever $x\in F$ nonzero.

Index of $Q \equiv \text{Ind}Q$, is largest possible dim. of subspace on which $Q$ is negative definite.








cf. 1.1. Morse's Lemma of  Ch. 6, pp. 145, Morse Theory of Hirsch (1997) \cite{dgdt:MHirsch1997}

\begin{lemma}[Morse's Lemma]
  Let $p\in M$ be nondegenerate critical pt. of index $k$ of $C^{r+2}$ map $f:M\to \mathbb{R}$, $1\leq r \leq \omega$.

  Then $\exists \, C^r$ chart $(\varphi,U)$ at $p$ s.t.
  \begin{equation}
    f\varphi^{-1}(u_1 \dots u_n) = f(p) - \sum_{i=1}^k u_i^2 + \sum_{i = k+1}^n u_i^2
  \end{equation}
  \end{lemma}

Let ${\,}^TQ \equiv Q^T$ denote tranpose of matrix $Q$.

\begin{lemma}
  Let $A = \text{diag}\lbrace a_1 , \dots , a_n \rbrace$ diagonal $n\times n$ matrix, with diagonal entries $\pm 1$.

  Then $\exists \, $ neighborhood $N$ of $A$ in vector space of symmetric $n\times n$ matrices, $C^{\infty}$ map
  \begin{equation}
P:N \to GL(n,\mathbb{R})
  \end{equation}
  s.t. $P(A)=I$, and if $P(B) = Q$, then $Q^T BQ = A$
\end{lemma}

\begin{proof}
  Let $B = [b_{ij}]$ be symmetri matrix near $A$ s.t. $b{11} \neq 0$ and $b_{11}$ has same sign as $a_1$.

    Consider $x=Ty$ where
    \[
\begin{aligned}
&  x_1 = \left[ y_1 - \frac{b_{12}}{b_{11}} y_2 - \dots - \frac{b_{1n}}{b_{11}} y_n \right] / \sqrt{ |b_n | } \\ 
& x_k = y_k \text{ for } k = 2, \dots n 
  \end{aligned}
    \]
  \end{proof}


\section{Lagrange multipliers}

From \emph{wikipedia:Lagrange multiplier}, \url{https://en.wikipedia.org/wiki/Lagrange_multiplier},
find local minima (maxima), pt. $a\in N$, s.t. $\exists \, $ neighborhood $U$ s.t. $f(x) \geq f(a)$ ($f(x) \leq f(a)$) \, $\forall \, x\in U$.

For $f:U\to \mathbb{R}$, open $U\subset \mathbb{R}^n$, find $x\in U$ s.t. $D_xf \equiv Df(x) =0$, check if Hessian $H_x f<0$.

Maxima may not exit since $U$ open.


References:

\href{http://www.math.uni.wroc.pl/~karch/analiza_nieliniowa/18_mnozniki-lagrangea.pdf}{Relative Extrema and Lagrange Multipliers}

Other interesting links:

\href{http://oai.cwi.nl/oai/asset/2552/2552A.pdf}{The Lagrange Multiplier Rule on Manifolds and Optimal Control of nonlinear systems}

